\documentclass[10pt,a4paper]{amsart}
\usepackage[margin=19mm]{geometry}
\usepackage{amsmath,amssymb,mathtools}
\usepackage[scr=rsfs,cal=euler]{mathalfa}
\usepackage{xfrac}
\usepackage[unicode,hidelinks,bookmarks=false]{hyperref}
\newcommand{\ie}{i.e.\ }
\newcommand{\cf}{cf.\ }
\newcommand{\resp}{respectively\ }
\newcommand{\Subsection}{\subsection}
\providecommand{\nobreakdash}{\nobreak\hbox{-}}
\setlength{\emergencystretch}{2em}
\multlinegap0pt
\usepackage{amssymb}
\usepackage{mathtools}
\usepackage{dsfont}
\usepackage{xcolor}
\usepackage{tikz}
\usepackage{bm}
\usepackage{cancel}
\usepackage{csquotes}
\newtheorem{lemma}{Lemma}
\DeclareMathOperator{\Cat}{Cat}
\def\One{\bm{1}}
\title{Tree-indexed sums of Catalan numbers}
\author{Alin Bostan, Valentin F\'eray, Paul Th\'evenin}
\date{Source: arXiv:2507.23557v2 (CC BY 4.0)}
\begin{document}
\maketitle

\noindent\emph{Source and licence.} Tree-indexed sums of Catalan numbers. Version arXiv:2507.23557v2 (CC BY 4.0), distributed by arXiv under the Creative Commons Attribution 4.0 International licence, http://creativecommons.org/licenses/by/4.0/, as recorded in the arXiv licence metadata for this exact version and shown on the abstract page https://arxiv.org/abs/2507.23557v2. Full licence notice: \texttt{licenses/arxiv-2507.23557-CC-BY-4.0.txt}.\par\smallskip

\noindent\textit{Editorial scope.} Counted unit: the article's lemma lem:twin\_varnothing (source0.tex lines 1274--1277). The article's complete original proof of it is delivered with it (source0.tex lines 1279--1314, one \texttt{proof} environment of 36 lines). No mathematics is written, repaired or re-derived: the statement and the proof are the article's own lines, in the article's own notation, and no retained line is substituted or re-flowed.  No further source-local context is needed: every cross-reference of every retained line resolves inside the two delivered spans.  One declared adaptation: exactly 1 ancillary strings inside the retained spans are omitted (image-inclusion commands and, where the source repeats a label, the repeated label definitions) and no other line differs from the source; the figure environments, with their placement, captions and labels, are kept, and the package records the omitted strings in fidelity receipts bound to the affected bodies.  No retained line contains a citation, so the wrapper carries no bibliography and no bibliography entry.  The retained lines contain the article's own diagram code, which the wrapper typesets with the standard tikz packages; no image file is delivered and the package declares no asset dependency.  Editorial formatting only, in addition: an AMS \texttt{amsart} preamble that restates the article's own declarations and macros used by the retained lines, namely the environments \texttt{lemma} on separate counters, together with the standard packages those lines need.  The retained environments print in this document as Lemma 1 in order of appearance, whereas the article prints the same items with the numbering of its own full document.  The complete native source of the article as distributed in the arXiv e-print for this exact version is bundled unchanged as \texttt{sources/182484/source0.tex}.
\begin{lemma}
\label{lem:twin_varnothing}
  $\begin{array}{c}  \end{array}$.   \\
    The same holds replacing white vertices with black vertices and the shift $+1$ with $-1$.\end{lemma}

\begin{proof}
Let $T$ be the tree on the left-hand side of Lemma~\ref{lem:twin_varnothing},
$w_1, w_2$ be the two white siblings with label $\varnothing$, 
and $p$ their common parent. 
Since $w_1$ and $w_2$ are decorated with $\varnothing$, the conditions $C_{w_1}$ and $C_{w_2}$ are void.
Furthermore, if $v$ is an ancestor of $p$, possibly $p$ itself, then the condition $C_{v}$ 
involves the sum $\ell_{w_1}+\ell_{w_2}$. Finally, if $v \ne w_1, w_2$ and $v$ is not an ancestor of $p$ (which we write $v\not\ge_T p$),
then the condition $C_{v}$ does not involve any of $\ell_{w_1}$, $\ell_{w_2}$.
Therefore the sum $S(T)$ can be written as follows:
\begin{multline*}
  S ( T ) = \sum_{(\ell_w)_{w\ne w_1,w_2} \atop (m_b)} \left[ \, \left( \prod_{w: w \ne w_1,w_2} \Cat_{\ell_w} t^{\ell_w} \right) 
\, \left( \prod_{b} \Cat_{m_b} t^{m_b} \right) \, \left( \prod_{v \not\ge_T p} \One[C_v] \right)\right. \\
\left. \cdot \sum_{L \ge 1}\,  \left( \prod_{v \ge_T p} \One[\tilde C_v] \right)
\, \left( \sum_{\ell_{w_1}+\ell_{w_2}=L-1} \Cat_{\ell_{w_1}} t^{\ell_{w_1}}\Cat_{\ell_{w_2}} t^{\ell_{w_2}}\right)\right],
\end{multline*}
where $\tilde C_v$ is the condition $C_v$ with the sum $\ell_{w_1}+\ell_{w_2}$ replaced by $L-1$.
The last sum in the above display simplifies as $t^{-L+1} \Cat_{L}$.
Having a sum over $L \ge 1$ is inconvenient, so we rewrite it as a sum over $L \ge 0$,
from which we remove the term corresponding to $L=0$. This gives
\begin{multline*}
  \hspace{-.5cm} S( T ) = t^{-1} \sum_{L,(\ell_w)_{w\ne w_1,w_2} \atop (m_b)}
 \Cat_{L} t^{L} \left( \prod_{w: w \ne w_1,w_2} \Cat_{\ell_w} t^{\ell_w} \right) 
\, \left( \prod_{b} \Cat_{m_b} t^{m_b} \right) \left( \prod_{v \not\ge_T p} \One[C_v] \right)\left( \prod_{v \ge_T p} \One[\tilde C_v] \right) \\
- t^{-1} \sum_{(\ell_w)_{w\ne w_1,w_2} \atop (m_b)}
\left( \prod_{w: w \ne w_1,w_2} \Cat_{\ell_w} t^{\ell_w} \right) 
\, \left( \prod_{b} \Cat_{m_b} t^{m_b} \right) \left( \prod_{v \not\ge_T p} \One[C_v] \right) \left( \prod_{v \ge_T p} \One[\tilde C_v] \right) ,
\end{multline*}
where all summation variables run over nonnegative integers.
The first (resp.~second) sum can be identified to $S(T_1)$,
where $T_1$ is the middle tree in Lemma~\ref{lem:twin_varnothing}
(resp.~$S(T_2)$,
where $T_2$ is the last tree in Lemma~\ref{lem:twin_varnothing}).
This proves the equation in the lemma.
The proof of the analogous statement with black vertices and shift $-1$
is similar.
\end{proof}

\end{document}
